% answers_short.tex —— 2014 年·重庆卷·理科 · 简短答案速查
% 由 tools/make_short_answers.py 自 answers.tex 的 \gksolhead/\gkanswer 提取，
% 并与 2014/2014_数学真题_重庆理_答案要点.md 逐题比对后生成，勿手改（改源 MD 后重跑）。

\begin{gknotice}
\textbf{使用说明}：本节只给最终答案便于快速核对，完整推导见解析版。答案为编写者独立推导，\textbf{非官方标准答案}。
\end{gknotice}

\gkpart{选择题}
\begin{center}\begin{tabular}{ccccccccccc}
\toprule
题号 & 1 & 2 & 3 & 4 & 5 & 6 & 7 & 8 & 9 & 10 \\
\midrule
答案 & A & D & A & C & C & D & B & B & B & A \\
\bottomrule
\end{tabular}\end{center}

\gkpart{填空题}
\begin{itemize}[leftmargin=2.6em]
  \item[\textbf{11.}] $\{7,9\}$。$\complement_U A=\{4,6,7,9,10\}$，与 $B=\{1,3,5,7,9\}$ 取交集得 $\{7,9\}$。
  \item[\textbf{12.}] $-\dfrac14$。令 $t=\log_2x$，则 $f=\dfrac t2\cdot2(1+t)=t^2+t=\left(t+\dfrac12\right)^2-\dfrac14$。
  \item[\textbf{13.}] $4\pm\sqrt{15}$。等边三角形边长等于半径 $2$，圆心到直线距离为 $\sqrt3$，即 $\dfrac{|2a-2|}{\sqrt{a^2+1}}=\sqrt3$，解得 $a^2-8a+1=0$。
  \item[\textbf{14.}] $4$。切割线定理 $PB\cdot PC=PA^2=36$ 且 $PC=PB+9$ 得 $PB=3$，$PC=12$；由 $\triangle PAB\sim\triangle PCA$ 得 $AB=\dfrac{PA\cdot AC}{PC}=4$。
  \item[\textbf{15.}] $\sqrt5$。直线为 $y=x+1$，曲线为 $y^2=4x$，联立得唯一公共点 $(1,2)$，故 $\rho=\sqrt{1^2+2^2}=\sqrt5$。
  \item[\textbf{16.}] $\left[-1,\dfrac12\right]$。$|2x-1|+|x+2|$ 的最小值为 $\dfrac52$，故 $a^2+\dfrac12a+2\leqslant\dfrac52$，即 $2a^2+a-1\leqslant0$。
\end{itemize}

\gkpart{解答题}
\begin{itemize}[leftmargin=2.6em]
  \item[\textbf{17.}] (1) $\omega=2$，$\varphi=-\dfrac\pi6$。相邻最高点距离为周期 $\pi$ 得 $\omega=2$；对称轴过最高点给出 $\dfrac{2\pi}3+\varphi=\dfrac\pi2+2k\pi$，结合 $\varphi$ 的范围得 $\varphi=-\dfrac\pi6$；(2) $\dfrac{\sqrt3+\sqrt{15}}8$。由 $\sin\left(\alpha-\dfrac\pi6\right)=\dfrac14$ 及 $\alpha-\dfrac\pi6\in\left(0,\dfrac\pi2\right)$ 得 $\cos\left(\alpha-\dfrac\pi6\right)=\dfrac{\sqrt{15}}4$，而 $\cos\left(\alpha+\dfrac{3\pi}2\right)=\sin\alpha$，按和角公式展开即得
  \item[\textbf{18.}] (1) $\dfrac5{84}$。三张全为 $1$ 有 $\mathrm C_4^3=4$ 种，全为 $2$ 有 $1$ 种，总取法 $\mathrm C_9^3=84$；(2) 分布列为 $P(X=1)=\dfrac{17}{42}$，$P(X=2)=\dfrac{43}{84}$，$P(X=3)=\dfrac1{12}$；$E(X)=\dfrac{47}{28}$。按三类数字的个数分类计数，$34+43+7=84$
  \item[\textbf{19.}] (1) $\dfrac{\sqrt3}2$。以 $O$ 为原点、菱形对角线为坐标轴建系，由 $\overrightarrow{MP}\cdot\overrightarrow{AP}=0$ 得 $PO^2=\dfrac34$；(2) $\dfrac{\sqrt{10}}5$。两半平面的法向量可取 $(\sqrt3,5,-2\sqrt3)$ 与 $(\sqrt3,-3,2\sqrt3)$，夹角余弦为 $\dfrac{\sqrt{15}}5$，故正弦值为 $\dfrac{\sqrt{10}}5$
  \item[\textbf{20.}] (1) $a=b=1$。$f'(x)=2ae^{2x}+2be^{-2x}-c$ 为偶函数给出 $a=b$，再由 $f'(0)=4a-c=4-c$ 得 $a=1$；(2) $f(x)$ 在 $\mathbf R$ 上单调递增。$c=3$ 时 $f'(x)=2e^{2x}+2e^{-2x}-3\geqslant4-3=1>0$；(3) $c>4$。$f'(x)\geqslant4-c$，当且仅当 $c>4$ 时 $f'$ 有两个变号零点，$f$ 才既有极大值又有极小值；$c=4$ 时 $f$ 仍严格递增，无极值
  \item[\textbf{21.}] (1) $\dfrac{x^2}2+y^2=1$。由 $|DF_1|=\dfrac{b^2}a$ 及两条件得 $c=1$、$a=\sqrt2$、$b=1$；(2) $\dfrac{4\sqrt2}3$。两交点关于 $y$ 轴对称，设为 $(\pm x_0,y_0)$；切线互相垂直给出 $x_0^2=(y_0-m)^2$，切线过焦点给出 $x_0-y_0=1$，联立椭圆方程得 $x_0=\dfrac43$，$y_0=\dfrac13$，$m=\dfrac53$，$r^2=\dfrac{32}9$
  \item[\textbf{22.}] (1) $a_2=2$，$a_3=1+\sqrt2$，$a_n=1+\sqrt{n-1}$。令 $b_n=a_n-1$ 得 $b_{n+1}=\sqrt{b_n^2+1}$，$b_1=0$，故 $b_n=\sqrt{n-1}$；(2) 存在，$c=\dfrac14$。令 $g(t)=\sqrt{(t-1)^2+1}-1$，由 $g$ 在 $(-\infty,1]$ 上递减、$g\left(\dfrac14\right)=\dfrac14$，归纳得 $0\leqslant a_{2n}<\dfrac14<a_{2n+1}$
\end{itemize}
